\section*{Code and Data Availability}

All simulations and analyses were performed using the open-source Python package \texttt{spinalmodes}, available at \url{https://github.com/sayuj1989-ops/life} (public fork of the working repository; the audit-fix branch underlying this revision is \texttt{fix/audit-2026-08-06}). Submission-time snapshots are preserved as git tags (\texttt{v1.0.0-submission}, \texttt{v1.4.7-editorial}, \texttt{v1.5-scaling-falsification}, \texttt{v1.7-consistency}); a Zenodo archival DOI is minted from the corresponding GitHub release (metadata in \texttt{.zenodo.json}). The computational environment, including dependencies such as \texttt{PyElastica}, \texttt{JAX}, \texttt{NumPy}, and optional \texttt{warp-lang} for GPU screening experiments, is specified in \texttt{pyproject.toml} and \texttt{requirements.txt}. All figure data are stored as CSV files and can be regenerated from the wrapper scripts in \texttt{scripts/experiments/}. The NVIDIA Warp screening outputs are stored under \texttt{results/nvidia\_warp\_beam\_sweep/} and are reproducible from \texttt{scripts/experiments/experiment\_nvidia\_warp\_beam\_sweep.py}. AlphaFold structural metrics were computed from AlphaFold Protein Structure Database snapshots cached in \texttt{research/alphafold\_countercurvature/data/} (manifest in \texttt{manifest.csv}); the manuscript protein tables are generated directly from the v6 metrics snapshot (\texttt{research/alphafold\_v6\_analysis/protein\_metrics.json}) by \texttt{scripts/analysis/regenerate\_table\_thermodynamic.py} and \texttt{regenerate\_table\_dissipation.py}, with the snapshot reconciliation printed by \texttt{scripts/analysis/protein\_snapshot\_provenance.py}. \textbf{No human-subject data, patient records, or proprietary clinical datasets were used. All correlations reported in this manuscript are model-internal} (simulation outputs analyzed against simulation parameters or against publicly available comparative-anatomy data); prospective external validation against AIS clinical cohorts is identified as future work.
